%% Lecture 19 -- Diamagnetism and paramagnetism: Larmor's theorem, the Langevin formula, and the Curie law
\section{Lecture 19 -- Diamagnetism and Paramagnetism (Langevin Theory)}\label{sec:lec19}

\begin{center}
\begin{tabular}{ll}
  \textbf{Date:}\quad 09/06/2026 & \\
\end{tabular}
\end{center}
\hrule\vspace{1.5em}

\subsection*{Overview}

Lecture 18 closed by sketching the microscopic origin of diamagnetism: switching on a
field $\bB$ makes every atomic electron precess, picking up a velocity
$\Delta\bvec{v}_n=\boldsymbol{\omega}\times\br_n$ at the Larmor frequency
$\boldsymbol{\omega}=-\tfrac{e}{2m_ec}\bB$, and the induced moment opposes $\bB$. That
derivation, however, leaned on the atom \emph{having} an orbital angular momentum
$\bvec{L}$ to precess. This lecture does three things:
\begin{itemize}
  \item \textbf{Frees Larmor's theorem from the assumption $\bvec{L}\neq0$:} by passing to
        a frame rotating at $\boldsymbol{\omega}$ we show the Coriolis force \emph{exactly
        cancels} the Lorentz force, so every electron acquires
        $\Delta\bvec{v}=\boldsymbol{\omega}\times\br$ whether or not the atom carries
        angular momentum. This is why diamagnetism is universal -- present in
        \emph{every} material.
  \item \textbf{Completes the Langevin diamagnetism formula}
        $\chi_{\rm dia}=-\tfrac{ne^2Z}{6mc^2}\langle r^2\rangle$ by averaging the induced
        moment over the precession, and estimates it numerically
        ($\chi_{\rm dia}\sim-10^{-9}$), with a digression on the \emph{classical electron
        radius} $r_c=e^2/mc^2$ that appears in the estimate.
  \item \textbf{Introduces paramagnetism:} for atoms with a \emph{permanent} moment
        (orbital or spin), precession alone produces no net alignment; it is
        \emph{thermal} agitation (collisions) that, through the Boltzmann distribution,
        biases the moments along $\bB$. This yields the \emph{Langevin function}
        $L(a)=\coth a-1/a$ and, in the weak-field limit, the \emph{Curie law}
        $\chi_{\rm para}=n\mu^2/3T>0$, some two orders of magnitude larger than the
        diamagnetic response.
\end{itemize}

%------------------------------------------------------------
\subsection{Larmor's Theorem Without Assuming $\bvec{L}$}\label{subsec:lec19-larmor}

\subsubsection*{Physical Motivation}
The previous derivation started from a precessing angular momentum,
$d\bvec{L}/dt=\boldsymbol{\omega}\times\bvec{L}$, and read off that each electron is carried
around with the precession, $\Delta\bvec{v}_n=\boldsymbol{\omega}\times\br_n$. But many
diamagnets (ground-state hydrogen, paired-electron molecules) have \emph{zero} total
angular momentum -- there is nothing to precess. We need to show the velocity increment
$\Delta\bvec{v}=\boldsymbol{\omega}\times\br$ is correct \emph{regardless} of $\bvec{L}$.
The cleanest argument never mentions $\bvec{L}$ at all: it works directly with the
equation of motion of a single electron, seen from a rotating frame.

\subsubsection*{The Rotating-Frame Argument}
Take one atomic electron undergoing some (arbitrary, classical) bound motion. Before the
field is switched on its equation of motion in the lab is
\begin{equation}\label{eq:lec19-eom-lab0}
  m\,\dot{\bvec{v}}=\bvec{F}_{\rm bind},
\end{equation}
with $\bvec{F}_{\rm bind}$ the (real) forces binding it to the nucleus. Switching on a
uniform $\bB$ adds the Lorentz force, so
\begin{equation}\label{eq:lec19-eom-lab}
  m\,\dot{\bvec{v}}=\bvec{F}_{\rm bind}+\frac{e}{c}\,\bvec{v}\times\bB .
\end{equation}
We want to understand the \emph{effect of the new term}. Pass to a frame rotating with
angular velocity
\begin{equation}\label{eq:lec19-omega}
  \boldsymbol{\omega}=-\frac{e}{2m_ec}\,\bB
\end{equation}
(the same $\boldsymbol{\omega}$ we can always build from $e,m,c$ and the applied $\bB$).
In a rotating frame two fictitious forces appear, the centrifugal force and the Coriolis
force,
\begin{equation}\label{eq:lec19-coriolis}
  \bvec{F}_{\rm cf}=-m\,\boldsymbol{\omega}\times(\boldsymbol{\omega}\times\br),
  \qquad
  \bvec{F}_{\rm Cor}=2m\,\bvec{v}'\times\boldsymbol{\omega},
\end{equation}
where $\bvec{v}'$ is the velocity \emph{measured in the rotating frame}. The centrifugal
term is $O(\omega^2)\propto B^2$; in the weak-field, non-relativistic regime of atomic
matter it is negligible and we drop it (see the pitfall note). For the Coriolis term use
the velocity-addition relation between lab and rotating frame,
\begin{equation}\label{eq:lec19-vprime}
  \bvec{v}'=\bvec{v}-\boldsymbol{\omega}\times\br,
\end{equation}
so that
\begin{equation}\label{eq:lec19-coriolis-expand}
  \bvec{F}_{\rm Cor}
  =2m\,(\bvec{v}-\boldsymbol{\omega}\times\br)\times\boldsymbol{\omega}
  =2m\,\bvec{v}\times\boldsymbol{\omega}
   \;\underbrace{-\,2m\,(\boldsymbol{\omega}\times\br)\times\boldsymbol{\omega}}_{O(\omega^2),\ \text{drop}} .
\end{equation}
Keeping the leading term and inserting \eqref{eq:lec19-omega},
\begin{equation}\label{eq:lec19-cancel}
  \bvec{F}_{\rm Cor}\simeq 2m\,\bvec{v}\times\boldsymbol{\omega}
  =2m\,\bvec{v}\times\!\left(-\frac{e}{2m c}\bB\right)
  =-\frac{e}{c}\,\bvec{v}\times\bB .
\end{equation}
This is precisely \emph{minus} the Lorentz force of \eqref{eq:lec19-eom-lab}. Therefore in
the rotating frame the two field-dependent terms annihilate and the equation of motion
collapses back to the field-free form
\begin{equation}\label{eq:lec19-eom-rot}
  m\,\dot{\bvec{v}}'=\bvec{F}_{\rm bind},
\end{equation}
identical to \eqref{eq:lec19-eom-lab0}.

\thm[thm:lec19-larmor]{Larmor's Theorem}{
  To leading order in $\bB$, the motion of the atomic electrons in a uniform field is the
  same motion they execute with no field, viewed from a frame rotating at the Larmor
  frequency $\boldsymbol{\omega}=-\tfrac{e}{2m_ec}\bB$. Equivalently, switching on $\bB$
  superposes on each electron the rigid rotation
  \begin{equation}\label{eq:lec19-larmor-boxed}
    \boxed{\Delta\bvec{v}_n=\boldsymbol{\omega}\times\br_n,
    \qquad
    \boldsymbol{\omega}=-\frac{e}{2m_ec}\,\bB,}
  \end{equation}
  \emph{independently} of whether the atom carries angular momentum.
}
\textit{Significance:} The induced circulation $\Delta\bvec{v}_n=\boldsymbol{\omega}\times\br_n$
is a property of the field and the charge-to-mass ratio alone, not of any pre-existing
orbit. Diamagnetism is therefore present in \emph{every} material; the bound currents it
needs are generated by the field itself.

\nt{Pitfall: the cancellation is \emph{exact only at order} $B$. The Coriolis force kills
the full Lorentz force, but the centrifugal force and the dropped term in
\eqref{eq:lec19-coriolis-expand} are $O(\omega^2)\propto B^2$. Keeping them would be
inconsistent -- we already discarded $B^2$ pieces elsewhere -- and physically they are
suppressed by $(\omega r/c)$-type factors that are tiny for atomic electrons
($v/c\lesssim10^{-2}$). ``Fictitious'' does not mean ``unreal'': in the rotating frame the
Coriolis force is as physical as gravity in an accelerated frame; here the frame change is
just a device to see what the added Lorentz force does.}

\subsubsection*{Method: turning ``switching on $\bB$'' into a kinematic rotation}
\begin{enumerate}
  \item Build the Larmor vector $\boldsymbol{\omega}=-\tfrac{e}{2m_ec}\bB$ from the applied
        field.
  \item Pass to the frame rotating at $\boldsymbol{\omega}$; write the two inertial forces
        \eqref{eq:lec19-coriolis}.
  \item Drop everything $O(\omega^2)$ (centrifugal, and the
        $\boldsymbol{\omega}\times\br$ piece of Coriolis).
  \item Check that the surviving Coriolis term equals $-\bvec{F}_{\rm Lorentz}$; conclude
        the rotating-frame motion is field-free.
  \item Transform back: the lab motion is the old motion plus the rigid rotation
        $\Delta\bvec{v}=\boldsymbol{\omega}\times\br$.
\end{enumerate}

%------------------------------------------------------------
\subsection{The Langevin Diamagnetism Formula}\label{subsec:lec19-langevin-dia}

\subsubsection*{Mathematical Setup}
The extra circulation \eqref{eq:lec19-larmor-boxed} is an extra current
$\Delta\Jvec=\sum_n e\,\Delta\bvec{v}_n\,\delta^{(3)}(\br-\br_n)$, hence an extra moment
$\Delta\bmu=\tfrac{1}{2c}\int d^3r\,\br\times\Delta\Jvec$. The $\delta$-functions collapse
the integral and, with $\Delta\bvec{v}_n=\boldsymbol{\omega}\times\br_n$,
\begin{equation}\label{eq:lec19-deltamu}
  \Delta\bmu=\frac{e}{2c}\sum_n \br_n\times(\boldsymbol{\omega}\times\br_n)
  =\frac{e}{2c}\sum_n\Bigl[\boldsymbol{\omega}\,r_n^2-\br_n(\boldsymbol{\omega}\cdot\br_n)\Bigr],
\end{equation}
using the BAC--CAB identity $\bvec{a}\times(\bvec{b}\times\bvec{a})=\bvec{b}\,a^2-\bvec{a}(\bvec{a}\cdot\bvec{b})$.

\subsubsection*{Projecting onto the Axes}
Choose $\bB=B\uvec{z}$, so $\boldsymbol{\omega}=\omega\uvec{z}$ with $\omega=-eB/2m c$.
Then $\boldsymbol{\omega}\cdot\br_n=\omega z_n$ and the three components of
\eqref{eq:lec19-deltamu} read
\begin{align}
  \Delta\mu_z &=\frac{e}{2c}\sum_n\bigl[\omega r_n^2-\omega z_n^2\bigr]
              =\frac{e\omega}{2c}\sum_n\bigl(x_n^2+y_n^2\bigr),\label{eq:lec19-muz}\\
  \Delta\mu_x &=-\frac{e\omega}{2c}\sum_n x_n z_n,\qquad
  \Delta\mu_y =-\frac{e\omega}{2c}\sum_n y_n z_n.\label{eq:lec19-muxy}
\end{align}

\subsubsection*{Averaging over the Precession}
We never observe instantaneous projections; the measured magnetization is a time average
over the (fast) precession. Averaging \eqref{eq:lec19-muxy}, the products $x_nz_n$ and
$y_nz_n$ are odd under the precession about $\uvec{z}$ and wash out,
\begin{equation}\label{eq:lec19-xy-vanish}
  \langle\Delta\mu_x\rangle=\langle\Delta\mu_y\rangle=0,
\end{equation}
leaving only the axial part. Summing over the $Z$ electrons of an atom and then over the
$n$ atoms per unit volume gives the magnetization (a density, hence the symbol $n$ for the
atomic number density),
\begin{equation}\label{eq:lec19-Mz-raw}
  M_z=\frac{n e\omega}{2c}\,Z\,\langle\rho^2\rangle,
  \qquad
  \rho^2\equiv x^2+y^2\ \ (\text{distance}^2\ \text{from the }z\text{-axis}).
\end{equation}
Insert $\omega=-eB/2mc$:
\begin{equation}\label{eq:lec19-Mz-rho}
  M_z=-\frac{n e^2 Z}{4 m c^2}\,\langle\rho^2\rangle\,B .
\end{equation}

\subsubsection*{Isotropy: from $\langle\rho^2\rangle$ to $\langle r^2\rangle$}
Before the field is applied a gas has no preferred direction, so
$\langle x^2\rangle=\langle y^2\rangle=\langle z^2\rangle=\tfrac13\langle r^2\rangle$ and
\begin{equation}\label{eq:lec19-rho-to-r}
  \langle\rho^2\rangle=\langle x^2+y^2\rangle=\tfrac23\langle r^2\rangle .
\end{equation}
(The field introduces a preferred axis only at $O(B)$, which would correct
\eqref{eq:lec19-Mz-rho} at $O(B^2)$ -- consistently negligible.) Hence
\begin{equation}\label{eq:lec19-Mz-final}
  M_z=-\frac{n e^2 Z}{6 m c^2}\,\langle r^2\rangle\,B .
\end{equation}

\dfn[dfn:lec19-langevin-dia]{Langevin Diamagnetic Susceptibility}{
  Because $|\chi|\ll1$ we may set $\Hvec\simeq\bB$ ($\mu\simeq1$) and read off
  $M=\chi H\simeq\chi B$:
  \begin{equation}\label{eq:lec19-chi-dia}
    \boxed{\chi_{\rm dia}=-\frac{n e^2 Z}{6 m c^2}\,\langle r^2\rangle<0 .}
  \end{equation}
  Here $n$ is the atomic number density, $Z$ the number of electrons per atom, and
  $\langle r^2\rangle$ the mean-square electron distance from the nucleus.
}
\textit{Significance:} The susceptibility is intrinsically \emph{negative} -- the induced
magnetization opposes $\bB$ -- and temperature-independent, the hallmark that
distinguishes Langevin diamagnetism from the Curie paramagnetism below. (Derived by
Langevin in 1905, the same year as special relativity.)

\nt{Pitfall: in \eqref{eq:lec19-chi-dia} ``$e$'' is the electron charge magnitude squared,
so the sign of $\chi$ comes entirely from the structure of the calculation, not from
$\mathrm{sign}(e)$. Also note $\langle r^2\rangle$ is the mean-square distance from the
\emph{nucleus}, while $\langle\rho^2\rangle$ in \eqref{eq:lec19-Mz-raw} is from the
\emph{axis}; the factor $2/3$ in \eqref{eq:lec19-rho-to-r} is exactly the bridge and is a
common place to drop a number.}

%------------------------------------------------------------
\subsection{Numerical Estimate and the Classical Electron Radius}\label{subsec:lec19-estimate}

\subsubsection*{Plugging in Numbers (CGS)}
To see how small $\chi_{\rm dia}$ is, take $Z\sim1$ (e.g.\ $\mathrm{H}_2$, He) and estimate
each factor:
\begin{itemize}
  \item \textbf{Number density} from the ideal-gas law $PV=Nk_BT\Rightarrow n=P/k_BT$. At
        room temperature ($T\sim300\,\mathrm{K}$) and one atmosphere,
        $n\sim2.7\times10^{19}\,\mathrm{cm^{-3}}$.
  \item \textbf{Atomic size} $\langle r^2\rangle^{1/2}\sim a_0\sim0.5\times10^{-8}\,\mathrm{cm}$
        (one \aa ngstr\"om).
  \item \textbf{The combination} $e^2/mc^2\equiv r_c=2.8\times10^{-13}\,\mathrm{cm}$, the
        \emph{classical electron radius}.
\end{itemize}
Writing $\chi_{\rm dia}=-\tfrac{nZ}{6}\,r_c\,\langle r^2\rangle$ and inserting these,
\begin{equation}\label{eq:lec19-chi-number}
  \boxed{|\chi_{\rm dia}|\sim 10^{-9}\ \ (\text{order of magnitude, and negative}).}
\end{equation}
\textit{Significance:} Even multiplied by $4\pi$ to form the permeability
$\mu=1+4\pi\chi$, this leaves $\mu$ indistinguishable from $1$ -- which is why ``ordinary''
matter looks non-magnetic. Paramagnetism, when present, is $\sim100\times$ stronger
($\chi\sim10^{-7}$), so it always masks the diamagnetic background.

\subsubsection*{Aside: Why $r_c=e^2/mc^2$ is Called the Classical Radius}
Model the electron as a charged sphere of radius $r_c$ with external field $E=e/r^2$. Its
electrostatic self-energy is
\begin{equation}\label{eq:lec19-selfenergy}
  U=\frac{1}{8\pi}\int_{r_c}^{\infty}E^2\,d^3r
   =\frac{1}{8\pi}\int_{r_c}^{\infty}\frac{e^2}{r^4}\,4\pi r^2\,dr
   =\frac{e^2}{2}\int_{r_c}^{\infty}\frac{dr}{r^2}=\frac{e^2}{2r_c},
\end{equation}
which \emph{diverges} as $r_c\to0$. Demanding that this electrostatic energy account for
the entire rest energy, $U\sim mc^2$, fixes
\begin{equation}\label{eq:lec19-rc}
  \boxed{r_c=\frac{e^2}{mc^2}\approx2.8\times10^{-13}\,\mathrm{cm}.}
\end{equation}
\textit{Significance:} $r_c$ is the scale at which a purely classical ``field energy
equals mass'' picture of the electron would break down. It is not a real size: the Compton
wavelength $\lambda_C=\hbar/mc\approx r_c/\alpha\sim100\,r_c$ (with the fine-structure
constant $\alpha=e^2/\hbar c\approx1/137$) is \emph{larger}, so quantum and relativistic
effects -- ultimately electron--positron pair creation -- intervene long before one probes
$r_c$. Asking ``how big is the electron'' is meaningless below $\lambda_C$, where any
measurement injects enough energy to create new particles.

%------------------------------------------------------------
\subsection{Paramagnetism: Why Precession Alone Does Nothing}\label{subsec:lec19-para-setup}

\subsubsection*{Physical Motivation}
Paramagnetism requires atoms with a \emph{permanent} magnetic moment $\bmu$, i.e.\ a
non-zero total angular momentum -- orbital $\bvec{L}$, electronic spin, or both (most
often spin, from an unpaired electron). A gas of such atoms has, before any field,
\emph{zero} net moment: the orientations are completely random, so $\sum\bmu=0$.

Switching on $\bB$ does \emph{not}, by itself, fix this. Every $\bmu$ feels the torque
$\bmu\times\bB$ and precesses, but they all precess at the \emph{same} Larmor rate
$\boldsymbol{\omega}$. A snapshot of the gas at any later time is just the initial random
configuration rigidly rotated, so the total moment stays zero. Precession reorients each
moment but creates no bias.

\subsubsection*{Conceptual Background: it is thermal, not mechanical}
What \emph{does} create a bias is the interaction of the atoms with each other -- in the
simplest picture, collisions. These let an atom's moment hop between orientations, and
because the orientation along $\bB$ has lower energy $U=-\bmu\cdot\bB$, the system relaxes
toward a thermal (Boltzmann) distribution that favours alignment. The competition between
this energetic preference and thermal randomization is the whole content of paramagnetism;
the microscopic details of the collisions drop out, which is exactly why one uses
\emph{statistical/thermal physics} rather than tracking individual dynamics.

\nt{This separates the two effects cleanly. Diamagnetism is a \emph{mechanical} response
(the field-induced current of Larmor's theorem) present in all matter. Paramagnetism is a
\emph{statistical} response (thermal alignment of pre-existing moments) present only when
the atoms carry angular momentum, and it overwhelms diamagnetism when it occurs.}

%------------------------------------------------------------
\subsection{The Langevin Function and the Curie Law}\label{subsec:lec19-curie}

\subsubsection*{Mathematical Setup: the Boltzmann Weight}
A moment at angle $\theta$ to $\bB=B\uvec{z}$ has energy
\begin{equation}\label{eq:lec19-U}
  U=-\bmu\cdot\bB=-\mu B\cos\theta .
\end{equation}
By the Boltzmann distribution the number of atoms per unit volume whose moment lies in the
angular strip $[\theta,\theta+d\theta]$ (integrating over the azimuth $\phi$, which nothing
depends on) is
\begin{equation}\label{eq:lec19-dn}
  dn=C\,e^{-\beta U}\sin\theta\,d\theta
    =C\,e^{a\cos\theta}\sin\theta\,d\theta,
  \qquad
  a\equiv\beta\mu B=\frac{\mu B}{T},
\end{equation}
with $\beta=1/k_BT$ and $k_B\equiv1$. The $\sin\theta\,d\theta$ is the solid-angle measure
(after the trivial $\phi$-integration); $C$ is a normalization.

\begin{figure}[h]
\centering
\begin{tikzpicture}[>=Stealth, thick, scale=1.0]
  % field axis
  \draw[->, blue!60!black, line width=1pt] (0,-1.6) -- (0,2.0) node[above] {$\bB=B\uvec{z}$};
  % the two cones bounding theta .. theta+dtheta
  \draw[gray] (0,0) -- (60:2);
  \draw[gray] (0,0) -- (75:2);
  % arc marking theta
  \draw[->] (0.8,0) arc (0:60:0.8);
  \node at (32:1.05) {$\theta$};
  % strip between the cones
  \draw[red!70!black, line width=1.2pt] (60:2) arc (60:75:2);
  \node[red!70!black, right] at (67:2.05) {$d\theta$};
  % a representative moment
  \draw[->, red!70!black, line width=1.2pt] (0,0) -- (67:1.5);
  \node[red!70!black] at (67:1.7) {$\bmu$};
\end{tikzpicture}
\caption{Geometry for the Boltzmann count. Moments pointing into the cone between $\theta$
and $\theta+d\theta$ have energy $-\mu B\cos\theta$; their number is
$dn\propto e^{a\cos\theta}\sin\theta\,d\theta$. The azimuth is integrated out by isotropy
about $\bB$.}
\label{fig:lec19-strip}
\end{figure}

\subsubsection*{Normalization}
Fixing $C$ by demanding $\int dn=n$ (the total number density),
\begin{equation}\label{eq:lec19-norm}
  n=C\int_0^{\pi}e^{a\cos\theta}\sin\theta\,d\theta
   =C\left[-\frac{1}{a}e^{a\cos\theta}\right]_0^{\pi}
   =C\,\frac{2\sinh a}{a}
  \ \Rightarrow\
  C=\frac{a\,n}{2\sinh a}.
\end{equation}

\subsubsection*{The Magnetization}
By azimuthal symmetry only the $z$-projection of each moment survives the sum; the
magnetization is
\begin{equation}\label{eq:lec19-Mz-int}
  M=M_z=\int \mu\cos\theta\,dn
   =\frac{\mu\,a\,n}{2\sinh a}\int_0^{\pi}\cos\theta\,e^{a\cos\theta}\sin\theta\,d\theta .
\end{equation}
Substituting $x=\cos\theta$ and differentiating the normalization integral under the
parameter $a$,
\begin{equation}\label{eq:lec19-Ideriv}
  \int_{-1}^{1}x\,e^{ax}\,dx
  =\frac{d}{da}\int_{-1}^{1}e^{ax}\,dx
  =\frac{d}{da}\!\left(\frac{2\sinh a}{a}\right)
  =\frac{2\,(a\cosh a-\sinh a)}{a^2}.
\end{equation}
Hence
\begin{equation}\label{eq:lec19-Mz-langevin}
  M=\frac{\mu\,a\,n}{2\sinh a}\cdot\frac{2(a\cosh a-\sinh a)}{a^2}
   =n\mu\left(\coth a-\frac1a\right).
\end{equation}

\thm[thm:lec19-langevin]{Langevin Magnetization}{
  A gas of $n$ moments $\mu$ per unit volume at temperature $T$ in field $B$ has
  magnetization along the field
  \begin{equation}\label{eq:lec19-langevin-boxed}
    \boxed{M=n\mu\,L(a),\qquad
    L(a)=\coth a-\frac1a,\qquad
    a=\frac{\mu B}{T},}
  \end{equation}
  where $L(a)$ is the \emph{Langevin function}.
}
\textit{Significance:} A single dimensionless ratio $a=\mu B/T$ -- the magnetic energy
measured in units of thermal energy -- controls the entire response, interpolating between
no alignment ($a\to0$) and complete saturation ($a\to\infty$).

\begin{figure}[h]
\centering
\begin{tikzpicture}[>=Stealth, scale=1.0]
  \draw[->] (0,0) -- (5.2,0) node[right] {$a=\mu B/T$};
  \draw[->] (0,0) -- (0,2.6) node[above] {$L(a)$};
  \draw[dashed] (0,2.0) -- (5.0,2.0) node[right] {$1$ (saturation)};
  % Langevin curve L(a)=coth a - 1/a, scaled so that L=1 -> y=2
  \draw[blue!70!black, line width=1pt]
    plot[domain=0.01:5,samples=60] (\x, {2*((cosh(\x)/sinh(\x)) - 1/\x)});
  % small-a tangent L ~ a/3  -> y = 2*a/3
  \draw[red!70!black, dashed] (0,0) -- (2.4,1.6) node[above left] {slope $\tfrac13$};
  \node[blue!70!black] at (3.6,1.55) {$L(a)$};
\end{tikzpicture}
\caption{The Langevin function $L(a)=\coth a-1/a$. For small $a$ it rises linearly with
slope $1/3$ (Curie regime); for large $a$ it saturates at $1$ (all moments aligned).}
\label{fig:lec19-langevin}
\end{figure}

\subsubsection*{Two Limits}
\paragraph{Saturation, $a\gg1$ (low $T$ or strong $B$).}
$\coth a\to1$ and $1/a\to0$, so
\begin{equation}\label{eq:lec19-sat}
  M\to n\mu :
\end{equation}
every moment is aligned, the maximum possible magnetization. Physically one cools a warm
gas (or raises $B$) until thermal agitation can no longer randomize the moments.

\paragraph{Curie regime, $a\ll1$ (the usual case).}
Expanding $L(a)=\tfrac{a}{3}-\tfrac{a^3}{45}+\dots$,
\begin{equation}\label{eq:lec19-curie-M}
  M\simeq n\mu\cdot\frac{a}{3}=\frac{n\mu^2 B}{3T}.
\end{equation}

\cor[cor:lec19-curie]{Curie's Law}{
  In the weak-field/high-temperature limit the paramagnetic susceptibility is
  \begin{equation}\label{eq:lec19-curie-boxed}
    \boxed{\chi_{\rm para}=\frac{n\mu^2}{3T}>0,}
  \end{equation}
  positive (magnetization \emph{along} $\bB$) and inversely proportional to temperature.
}
\textit{Significance:} The $1/T$ dependence is the experimental signature of paramagnetism
and the sharpest contrast with the temperature-independent diamagnetic
\eqref{eq:lec19-chi-dia}: warming the gas lets thermal motion fight the alignment, weakening
the response. With $\mu$ of order a Bohr magneton, $\chi_{\rm para}\sim10^{-7}$ at room
temperature -- about $100\times$ the diamagnetic background, which it therefore dominates.

\nt{Pitfall: $a=\mu B/T$ is genuinely \emph{tiny} in ordinary conditions
($\mu B\ll k_BT$), so the linear Curie law -- not the full Langevin function -- is what one
uses; saturation is a low-temperature / high-field laboratory regime. Keep $k_B$ explicit
($a=\mu B/k_BT$) if you are tracking units; here $k_B=1$.}

\subsubsection*{Method: computing a thermal magnetization}
\begin{enumerate}
  \item Write the orientation energy $U=-\bmu\cdot\bB=-\mu B\cos\theta$.
  \item Weight orientations by Boltzmann, $dn\propto e^{-\beta U}\sin\theta\,d\theta$, and
        fix the constant by $\int dn=n$.
  \item Average the field-projection: $M=\int\mu\cos\theta\,dn$; evaluate via $x=\cos\theta$
        and ``differentiate the partition integral under $a$''.
  \item Read off $M=n\mu\,L(a)$; expand $L(a)\approx a/3$ for the susceptibility, or take
        $L\to1$ for saturation.
\end{enumerate}

%------------------------------------------------------------
\subsection{Summary}\label{subsec:lec19-summary}

\begin{itemize}
  \item \textbf{Larmor's theorem} holds without assuming $\bvec{L}$: in the frame rotating
        at $\boldsymbol{\omega}=-\tfrac{e}{2m_ec}\bB$ the Coriolis force cancels the Lorentz
        force ($2m\bvec{v}\times\boldsymbol{\omega}=-\tfrac{e}{c}\bvec{v}\times\bB$), so each
        electron simply acquires $\Delta\bvec{v}=\boldsymbol{\omega}\times\br$. Hence
        diamagnetism is universal.
  \item \textbf{Langevin diamagnetism:} averaging the induced moment leaves only the axial
        part, giving $\chi_{\rm dia}=-\tfrac{ne^2Z}{6mc^2}\langle r^2\rangle<0$,
        temperature-independent and $\sim-10^{-9}$ numerically.
  \item The estimate uses the \textbf{classical electron radius}
        $r_c=e^2/mc^2\approx2.8\times10^{-13}\,$cm, obtained by equating the electrostatic
        self-energy $e^2/2r_c$ to $mc^2$; it is unphysical below the Compton wavelength
        $\lambda_C=\hbar/mc\sim100\,r_c$.
  \item \textbf{Paramagnetism} needs permanent moments and is thermal, not mechanical:
        precession alone keeps $\sum\bmu=0$; collisions drive Boltzmann alignment.
  \item The Boltzmann average gives the \textbf{Langevin function}
        $M=n\mu\,(\coth a-1/a)$, $a=\mu B/T$: saturation $M\to n\mu$ for $a\gg1$, and the
        \textbf{Curie law} $\chi_{\rm para}=n\mu^2/3T>0$ for $a\ll1$ -- positive, $\propto
        1/T$, and $\sim100\times$ stronger than diamagnetism.
\end{itemize}
